%%%%%%%%%%%%%%%%%%%%%%%%%%%%%%%%%%%%%%%%%
% Afstudeer scriptie
%%%%%%%%%%%%%%%%%%%%%%%%%%%%%%%%%%%%%%%%%

%----------------------------------------------------------------------------------------
%	PACKAGES AND OTHER DOCUMENT CONFIGURATIONS
%----------------------------------------------------------------------------------------

\documentclass[a4paper,12pt,oneside]{report} 
% Default font size is 12pt, it can be changed here

\usepackage[utf8]{inputenc}
\usepackage[T1]{fontenc}
\usepackage{palatino}

\usepackage{listings}

\usepackage{blindtext}
\usepackage[ampersand]{easylist}
\usepackage{fancyhdr}

\usepackage[xindy,toc]{glossaries}
\usepackage[font={small,it}]{caption}
\usepackage{tocloft}
%\usepackage[options]{natbib}
\usepackage[numbers]{natbib} 

\usepackage{hyperref}
\hypersetup{
    colorlinks,
    citecolor=black,
    filecolor=black,
    linkcolor=black,
    urlcolor=black
}

\usepackage[dutch]{babel}
\selectlanguage{dutch}

\usepackage[left=1.7in]{geometry} % Required to change the page size to A4
\geometry{a4paper} % Set the page size to be A4 as opposed to the default US Letter

\usepackage{graphicx} % Required for including pictures

\usepackage{float} % Allows putting an [H] in \begin{figure} to specify the exact location of the figure
\usepackage{wrapfig} % Allows in-line images such as the example fish picture

\usepackage{lipsum} % Used for inserting dummy 'Lorem ipsum' text into the template

\linespread{1.2} % Line spacing



\graphicspath{{./img/}} % Specifies the directory where pictures are stored


\usepackage{listings}
\usepackage[usenames,dvipsnames]{xcolor}
\usepackage{textcomp}

\definecolor{javared}{rgb}{0.6,0,0} % for strings
\definecolor{javagreen}{rgb}{0.25,0.5,0.35} % comments
\definecolor{javapurple}{rgb}{0.5,0,0.35} % keywords
\definecolor{javadocblue}{rgb}{0.25,0.35,0.75} % javadoc
\definecolor{lbcolor}{rgb}{0.9,0.9,0.9} % javadoc
\definecolor{darkblue}{rgb}{0.0,0.0,0.6}
\definecolor{cyan}{rgb}{0.0,0.6,0.6}
\definecolor{code-block}{rgb}{0.97,0.97,0.97}
 
\lstset{
	language=Java,
	basicstyle=\scriptsize,
	keywordstyle=\color{javapurple}\bfseries,
	stringstyle=\color{javared},
	commentstyle=\color{javagreen},
	morecomment=[s][\color{javadocblue}]{/**}{*/},
	backgroundcolor=\color{lbcolor},
	numbers=left,
	numberstyle=\color{black},
	stepnumber=1,
	numbersep=10pt,
	tabsize=2,
	showspaces=false,
	showstringspaces=false
	breaklines=true
}



\usepackage[compact]{titlesec}
\usepackage{pdfpages}

\setlength\aftertitleunit{\baselineskip}



\fancypagestyle{plain}{}


\newlength\titleindent
\setlength\titleindent{2cm}

\pretocmd{\paragraph}{\stepcounter{subsection}}{}{}
\pretocmd{\subparagraph}{\stepcounter{subsubsection}}{}{}

\titleformat{\chapter}[block]
  {\normalfont\huge\bfseries}{\llap{\parbox{\titleindent}{\thechapter\hfill}}}{0pt}{}

\titleformat{\section}
  {\normalfont\Large\bfseries}{\llap{\parbox{\titleindent}{\thesection\hfill}}}{0em}{}

\titleformat{\subsection}
  {\normalfont\large}{\llap{\parbox{\titleindent}{}}}{0em}{\bfseries}

\titleformat{\subsubsection}
  {\normalfont\normalsize}{\llap{\parbox{\titleindent}{\thesubsubsection}}}{0em}{\bfseries}

\titleformat{\paragraph}[runin]
  {\normalfont\large}{\llap{\parbox{\titleindent}{\thesubsection\hfill}}}{0em}{}

\titleformat{\subparagraph}[runin]
  {\normalfont\normalsize}{\llap{\parbox{\titleindent}{\thesubsubsection\hfill}}}{0em}{}

\titlespacing*{\chapter}{0pt}{0pt}{20pt}
\titlespacing*{\subsubsection}{0pt}{3.25ex plus 1ex minus .2ex}{1.5ex plus .2ex}
\titlespacing*{\paragraph}{0pt}{3.25ex plus 1ex minus .2ex}{0em}
\titlespacing*{\subparagraph}{0pt}{3.25ex plus 1ex minus .2ex}{0em}

\setcounter{tocdepth}{1}
\cftsetindents{chapter}{0in}{0.5in}
\cftsetindents{section}{0in}{0.5in}

\lstdefinestyle{code-block}{
  language=bash,
  basicstyle=\small\ttfamily,
  linewidth=0.9\linewidth,
  breaklines=true
}

\lstset{
   backgroundcolor=\color{blue!20}
}

%----------------------------------------------------------------------------------------
%	Glossary
%----------------------------------------------------------------------------------------


\makeglossaries


%----------------------------------------------------------------------------------------
%	Achronyms
%----------------------------------------------------------------------------------------



\begin{document}

%----------------------------------------------------------------------------------------
%	Heading
%----------------------------------------------------------------------------------------

\pagestyle{fancy}
\renewcommand{\headrulewidth}{0.0pt}
\headheight = 54pt
 
\lhead{\includegraphics[width=0.4\linewidth]{finan_logo.eps}}
\rhead{\includegraphics[width=0.4\linewidth]{hanze_logo.png}}






%----------------------------------------------------------------------------------------
%	TITLE PAGE
%----------------------------------------------------------------------------------------

\begin{titlepage}
\oddsidemargin 1cm

\newcommand{\HRule}{\rule{\linewidth}{0.5mm}} % Dopenrightefines a new command for the horizontal lines, change thickness here

\center % Center everything on the page

\textsc{\LARGE Hanzehogeschool}\\[1.0cm] % Name of your university/college

\textsc{\large \textit{"Onder titel"} }\\[0.5cm] % Major heading such as course name


\HRule \\[0.4cm]
{ \huge \bfseries Afstudeer scriptie}\\[0.4cm] % Title
\HRule \\[6cm]



\includegraphics[width=\linewidth]{finan_logo.eps}\\
\today, Zwolle
\end{titlepage}

\large
\emph{Auteur:}\\
Erwin Oldenkamp\\
 \textsc{282472} \\
Opleiding: Informatica \\



\emph{Stagedocent:}\\
Henkjan Hekman\\
Hanzehogeschool Groningen \\



\emph{Stagebedrijf:}\\
    Topicus Finan\\
	Wouter van der Meulen Kunee\\
	Koggelaan 3A\\
	8017 JH, Zwolle


\pagenumbering{gobble}

%----------------------------------------------------------------------------------------
%	Woordvooraf
%----------------------------------------------------------------------------------------
\chapter*{Woord vooraf}
\lipsum[1]

%----------------------------------------------------------------------------------------
%	Samenvatting
%----------------------------------------------------------------------------------------
\chapter*{Samenvatting}
\lipsum[1]





%----------------------------------------------------------------------------------------
%	Index
%----------------------------------------------------------------------------------------
\newpage
\renewcommand*\contentsname{Inhoud}
\tableofcontents
\cleardoublepage 
\pagenumbering{arabic}




%----------------------------------------------------------------------------------------
%	Inleiding
%----------------------------------------------------------------------------------------

\chapter{Inleiding} 
Inleiding

%----------------------------------------------------------------------------------------
%	Topicus Finan BV
%----------------------------------------------------------------------------------------

\section{Topicus Finan BV} 
Topicus Finan BV is een organisatie die oplossingen biedt voor de ondersteuning van financiële analyses.
Finan is marktleider op het gebied van software voor financiële en niet-financiële verantwoordingsstromen (waaronder credit rating toepassingen) en biedt SaaS oplossingen op het gebied van klantspecifieke berichtenportalen en financiële front-, mid- en backoffice applicaties. 
Topicus Finan BV bestaat uit 3 afdelingen, Accountancy, Banking en Retail Banking. 

\subsection{Accountancy}
\lipsum[1]
\subsection{Banking}
\lipsum[1]
\subsection{Retail Banking}
Retail Banking houdt zich bezig met het automatiseren van financiële processen. 
Sinds 2008 is Wehkamp.nl een van de grootste klanten van de afdeling Retail Banking.
Het Finan Retail Banking systeem bevat een hoop informatie met betrekking tot de klanten van een leverancier van betalingsfaciliteiten. 
Het systeem bestaat uit twee applicaties, CAS en IRIS. 

CAS is een crediteuren en debiteuren administratie systeem. 
In dit systeem kan alles geautomatiseerd worden qua debiteuren en crediteuren.
IRIS is een geautomatiseerd debiteuren dossier. 
Met dit systeem kan de servicedesk klanten te woord staan die vragen of opmerkingen hebben over de financiële afwerking. 




\subsection{Organisatiestructuur}
\lipsum[2]

\subsection{De primaire processen}
\lipsum[3]


\section{Probleemstelling}
\lipsum[4]

\section{Inhoud van dit rapport}
\lipsum[5]

\chapter{Plan van aanpak}
\lipsum[6-7]

%----------------------------------------------------------------------------------------
%	Keuze platform
%----------------------------------------------------------------------------------------

\chapter{Keuze platform}
\lipsum[8]
\section{Criteria}
\lipsum[9]

\section{Uiteindelijke keuze}
\lipsum[10]



%----------------------------------------------------------------------------------------
%	CONCLUSION
%----------------------------------------------------------------------------------------

\chapter{Conclusie}
\lipsum[11-14]




\appendix
%----------------------------------------------------------------------------------------
%	Persoonlijke ontwikkeling
%----------------------------------------------------------------------------------------
\chapter{Persoonlijke ontwikkeling}
\lipsum[1]

%----------------------------------------------------------------------------------------
%	woordenlijst en bronnen
%----------------------------------------------------------------------------------------
\newpage

\printglossary


\renewcommand{\bibname}{Bronvermeldingen}
\bibliographystyle{plain}
\bibliography{bron}
\nocite{*}

%----------------------------------------------------------------------------------------



\end{document}
